% =====================================================================
\section{Controllable Objects}\label{sec:control}
The scene has \nControllables{} controllable objects, \nPossessable{} of them \emph{possessable}: nine machines, the UAV (manual flight), tank, helicopter, AA gun, radar, blast doors and gate barrier (this section) and the SAM and MLRS launchers (Section~\ref{sec:v5media}); since v6 also the \nSoldiers{} soldiers (\texttt{F3}, Section~\ref{sec:soldiers}) and the SKY clock that carries the sun (Section~\ref{sec:daycycle}). Windsock and flag animate but are not possessable. Each implements \code{Controllable} (\code{scene/Controllable.h}): every frame only the possessed object receives the sampled \code{ControlInput} (move, turn, lift, aim, fire, action) and writes the local transforms of its own node tree in \code{update()}. In walk mode \texttt{Enter} toggles a door or barrier within 3.5\,m (T128). Keys: Appendix~\ref{app:controls}.

\begin{table}[h]
\centering\footnotesize
\setlength{\tabcolsep}{3pt}
\caption{Controllable objects. DOF = independently animated degrees of freedom.}
\label{tab:ctl}
\begin{tabular}{@{}clp{1.7cm}p{5.2cm}p{5.4cm}@{}}
\toprule
Key & Object (HUD name) & Class & Degrees of freedom & Camera views\\\midrule
5 & MBT-1 TANK & \code{Tank} & position (x, z), yaw, left/right track, turret yaw, gun pitch, recoil, hull heave/pitch, 2 antennas, MG ring & chase (9\,m behind, 55$^\circ$), sight (28$^\circ$)\\
6 & MH-6 HELICOPTER & \code{Helicopter} & position (3), yaw, pitch, roll, rotor speed, 5 blades (flap, feather), tail rotor, light pan/tilt & chase / cockpit\\
7 & AA-1 GATLING & \code{AaGun} & mount yaw, cradle pitch, cluster spin, search antenna & chase / sight\\
8 & AIR-DEFENCE RADAR & \code{Radar} & azimuth, elevation tilt, mode, TX lamp & side view (13\,m away), operator view along the beam (60$^\circ$)\\
9 & BLAST DOORS & \code{BlastDoors} & two leaves (slide), beacon spinners & door camera at $(0,-0.3,9)$\\
0 & UAV-1 DRONE & \code{DroneControl} & position (3), yaw, pitch, roll, gimbal yaw/pitch, rail rounds & chase / gimbal camera (40$^\circ$)\\
-- (\texttt{[ ]}) & GATE BARRIER & \code{GateBarrier} & boom angle, lamps & near the gate\\
F3 & SOLDIER 1--8 (v6) & \code{Soldier} & position (x, z), heading, view pitch, 2 legs (hip, knee, ankle) by gait + IK, 2 arms (shoulder 3, elbow) by swing or rifle IK, spine, head, stance (stand / crouch / wave) & chase (3.2\,m behind), first person at eye height 1.62\,m (T212)\\
-- (\texttt{[ ]}) & SKY / TIME OF DAY (v6) & \code{SkyClock} & clock hour (moves the sun sphere), auto / manual, phase jump & from the camp towards the key light\\\bottomrule
\end{tabular}
\end{table}

% ---------------------------------------------------------------------
\begin{figure}[h]
\centering
\picsheet{sheet_controllables}{\linewidth}
\caption{\sheetcaption{sheet_controllables}}
\end{figure}
\subsection{Tank: arc-length track links, skid steer, stabilised turret}
\sheetfig{s14_tank}
\paragraph{Track-link arc-length parameterisation.} The closed track path is a stadium in the hull's $(z,y)$ plane around the
sprocket ($z=-2.8$) and the idler ($z=2.8$), centre height 0.55\,m, radius $R=0.5$\,m, with perimeter
$P=2(2.8-(-2.8))+2\pi R=14.342$\,m. Each side keeps one number, the travelled distance
$\phi_\text{belt}\mathrel{+}=v_\text{track}\,\Delta t$; link $i$ of 40 sits at arc length
$s_i=(\phi_\text{belt}+P\,i/40)\bmod P$, mapped by \code{tankpath::point(s)} to $(z,y,\varphi)$ (top run, front arc, bottom run, rear arc). The link translates by $(0,y,z)$ and
$R_x(\varphi)$ turns its local $+y$ onto the outward normal. The wheels spin by
$\phi_\text{belt}/r_\text{wheel}$, so tread and wheels agree. All 80 links stay on the path (T80) and the belt advances by $v\,\Delta t$ (T81: 16.0000\,m for 16\,m).

\paragraph{Skid steer.} Left and right track speeds are $v_{l,r}=f\pm1.2\,u$ with forward command $f$ (max 8\,m/s, 12\,m/s boosted, 3\,m/s reverse) and turn input $u$, approached at 3.5\,m/s$^2$ (7\,m/s$^2$ braking). The body speed is
$v=(v_l+v_r)/2$ and the yaw rate $\omega=(v_r-v_l)/2.2$ (track separation 2.2\,m; positive yaw is about $+Y$). Equal speeds drive straight (T82); $\pm1.2$\,m/s pivots at $-62.5^\circ$/s (T83). Collision is an
oriented-hull vs.\ box separating-axis test against the static solids; a blocked tank tries to pivot, then slide, then stops; it stays within $r\le31$\,m (T87).
The hull height is the mean terrain height under four track corners (T86: worst error 2.3\,mm).

\paragraph{Suspension, recoil and the stabilised turret.} The hull pitch is a damped spring
$\ddot x=\omega^2(x_\text{target}-x)-2\zeta\omega\dot x$ with $\omega=2\pi\cdot1.6$, $\zeta=0.35$, driven by
$x_\text{target}=\mathrm{clamp}(-0.012\,a,\pm0.06)$ (acceleration $a$): the tank squats when it accelerates and rocks back when it
fires. The antennas use the same spring ($\omega=2\pi\cdot2.2$, $\zeta=0.12$). The turret keeps its \emph{world} yaw when
the hull turns (T84b: hull $-114.8^\circ$, turret unchanged at $-40^\circ$); it slews at 40$^\circ$/s, the gun at 20$^\circ$/s over
$-8\ldots+20^\circ$ (T84). The 120\,m/s shell recoils the barrel 0.35\,m (back in 0.6\,s); reload 3.5\,s, 18 rounds (T85); the range check is in Section~\ref{sec:game} (T101).

% ---------------------------------------------------------------------
\subsection{Helicopter: cyclic flapping and a spring-damper attitude model}
\sheetfig{s14_helicopter}
\paragraph{State machine.} \emph{Parked, Startup, Hover, Fly, Landing.} \code{\textbackslash} starts the engine (rotor
spool-up 6\,s, T90; 300\,rpm, tail rotor 5$\times$ faster) or requests a landing, accepted only at ground speed $\le6$\,m/s, height $\le20$\,m and skid slope $\le0.15$
(T94: slope 0.31 refused, reason shown). Spool-down takes 8\,s; released in flight the helicopter hovers (T95).

\paragraph{Flight limits.} Cruise 15\,m/s (boost 25, reverse 7), climb 4\,m/s, 0.5--40\,m above ground, radius $\le60$\,m (T93), acceleration 8\,m/s$^2$ (12 boosted), yaw 55$^\circ$/s.

\paragraph{Spring-damper attitude.} Pitch and roll follow acceleration and turn rate through
$\ddot x=\omega^2(x_\text{target}-x)-2\zeta\omega\dot x$ with $\omega=5$\,s$^{-1}$, $\zeta=0.8$, bounded to $22^\circ$ pitch and $25^\circ$ roll
(T92: max pitch 22.2$^\circ$, roll 25.5$^\circ$ at 24.5\,m/s).

\paragraph{Rotor cyclic.} The rotor head carries a tip-path-plane node that tilts by the cyclic angles $(\tau_f,\tau_l)$, which
approach their targets exponentially ($T=0.2$\,s). Each of the five blades, $72^\circ$ apart (T91), sits under a flap node rotated by $R_x(-\beta)$, $R_z(\vartheta)$ with
$\beta=\beta_\text{cone}-\tfrac12\tau_f\cos\psi-\tfrac12\tau_l\sin\psi$ and the feather
$\vartheta=\vartheta_\text{coll}+0.6(\tau_f\sin\psi-\tau_l\cos\psi)$ ($\psi$ is the blade azimuth); the blades droop at low rpm. A
translucent blur disc fades in with rotor speed (T96: 0 $\to$ 0.26). Below 4\,m with the rotor above 60\,\% the helicopter kicks up dust (\code{Combat::emitDust}).
The arrows or mouse aim the searchlight (pan $\pm150^\circ$, tilt $-10\ldots90^\circ$ down).

% ---------------------------------------------------------------------
\subsection{AA Gatling gun: spin-up and tracers}
\sheetfig{s14_aagun}
The gun slews in azimuth at 60$^\circ$/s through 360$^\circ$ and in elevation from $-5^\circ$ to $+85^\circ$ (T89). While \texttt{Enter} is held the barrel
cluster spins up linearly to $5$\,rev/s in 0.5\,s and coasts down in 1.5\,s; rounds are released only at full spin, 30 rounds/s ($6$ barrels
$\times5$ rev/s), as tracers at 1000\,m/s with 3.5\,mrad dispersion, from 1550 rounds (T88: first round at
0.517\,s, 60 in 2\,s; spin 5.00 $\to$ 2.50 $\to$ 0.00\,rev/s). Casings are a pool of 20 ballistic parts. Unpossessed, the mount sweeps slowly and never fires (T89: moved 157$^\circ$, fired 0).
The search antenna turns at 1\,rev/s (3 when possessed); the sight camera sits on the cradle.

% ---------------------------------------------------------------------
\subsection{Radar: three modes}
\sheetfig{s14_radar}
\emph{ROTATE}: azimuth advances 72$^\circ$/s (12\,rpm, one turn in 5\,s, T97); \emph{SECTOR}: a sine sweep of $\pm45^\circ$ around a centre
that the user aims, period 6\,s (T98); \emph{STOP}: manual azimuth. Tilt is limited to $-5\ldots+20^\circ$
(T99). Released, the radar returns to ROTATE. Azimuth is measured clockwise from $+Z$, i.e.\ $-$yaw about $+Y$.

% ---------------------------------------------------------------------
\subsection{Blast doors: eased sliding leaves}
\sheetfig{s14_doors}
The two leaves (1.45\,m wide, six parts each, opening 2.6\,m) slide 1.5\,m into the wall pockets with
$x(t)=\pm1.5\,s(t)$, $s(u)=3u^2-2u^3$, over 4\,s. The states are Closed, Opening, Open, Closing, reversible mid-way
(T65: 4.02\,s one way). The doors \emph{start open}; two beacons spin at 300$^\circ$/s while they move. A closed door is a solid box in the walker's collision via
\code{dynamicObstacles()} (T64, T121), as is the lowered barrier arm.

% ---------------------------------------------------------------------
\subsection{Gate barrier, windsock and flag}
\sheetfig{s14_gate}
\textbf{Barrier.} The boom rotates $0\to85^\circ$ about its pivot in 2.5\,s with the same smoothstep easing
$\alpha=85^\circ s(t)$; its lamp blinks at 3\,Hz while moving; below $15^\circ$ the arm is an obstacle box (T73: half-way at 1.25\,s, tip height 6.04\,m). \textbf{Windsock.} A swivel follows the wind direction $200^\circ+25^\circ\sin(0.21t)+10^\circ\sin(0.57t+1.3)$; five chained
segments droop by $20^\circ(1-w)+2.5^\circ$ plus a flutter with a phase lag down the chain ($w\in[0.35,1]$ is the wind strength; T75: 16$^\circ$ at $w=0.36$, 0.7$^\circ$ at $w=1$). \textbf{Flag.} Twelve strips in a chain, each yawed by a travelling
sine wave $\mathrm{yaw}_i=A_i\sin(0.7\,i-\omega t)$, $A_i=(2+4w)(0.3+0.7\,i/11)$ and $\omega=2\pi(0.6+1.4w)$, plus a small pitch term; the rotations accumulate down the chain, so the tip distance stays 2.50--2.52\,m (T74): the flag neither stretches nor collapses.

% ---------------------------------------------------------------------
\subsection{UAV in manual flight}
\sheetfig{s14_uav}
Key \texttt{0} takes the first UAV: forward/back 6\,m/s (9 with Shift), yaw 70$^\circ$/s, climb 3\,m/s, acceleration $\le4$\,m/s$^2$, height 1--40\,m, radius $r\le34$\,m. Pitch and roll follow
$4\,\mathbf v\cdot\hat f+3\,\mathbf a\cdot\hat f$ (clamped to $\pm18^\circ$) with exponential smoothing. The arrows or mouse turn the gimbal (60$^\circ$/s yaw, 40$^\circ$/s pitch,
5--88$^\circ$); \texttt{\textbackslash} toggles \emph{gimbal assist}, in which the gimbal keeps following the tracked target. \texttt{Enter} releases a rail round at the ground point
under the camera centre (Section~\ref{sec:game}); autonomous squad flight is Section~\ref{sec:squad}.

\begin{figure}[h]
  \centering
  \begin{subfigure}{0.325\linewidth}\picfour{possess_tank_sight}{\linewidth}\caption{Tank sight (\texttt{5})}\end{subfigure}
  \begin{subfigure}{0.325\linewidth}\picfour{possess_heli_cockpit}{\linewidth}\caption{Helicopter cockpit (\texttt{6})}\end{subfigure}
  \begin{subfigure}{0.325\linewidth}\picfour{possess_aa_sight}{\linewidth}\caption{AA gun sight (\texttt{7})}\end{subfigure}\\[1mm]
  \begin{subfigure}{0.325\linewidth}\picfour{possess_uav_manual}{\linewidth}\caption{UAV, manual (\texttt{0})}\end{subfigure}
  \begin{subfigure}{0.325\linewidth}\picfour{possess_doors}{\linewidth}\caption{Blast doors (\texttt{9})}\end{subfigure}
  \begin{subfigure}{0.325\linewidth}\picfour{heli_takeoff}{\linewidth}\caption{Helicopter take-off}\end{subfigure}
  \caption{Possession views with the HUD (captures from \code{Capture.cpp}).}
  \label{fig:possess}
\end{figure}

\begin{figure}[h]
  \centering
  \begin{subfigure}{0.325\linewidth}\picfour{tank_closeup}{\linewidth}\caption{Tank at the motor pool}\end{subfigure}
  \begin{subfigure}{0.325\linewidth}\picfour{tank_fire}{\linewidth}\caption{Tank firing}\end{subfigure}\\[1mm]
  \begin{subfigure}{0.325\linewidth}\picfour{helipad}{\linewidth}\caption{Helipad and helicopter}\end{subfigure}
  \begin{subfigure}{0.325\linewidth}\picfour{radar}{\linewidth}\caption{Radar}\end{subfigure}
  \begin{subfigure}{0.325\linewidth}\picfour{aa_gun_burst}{\linewidth}\caption{AA gun burst}\end{subfigure}
  \caption{The machines in the scene.}
  \label{fig:machines}
\end{figure}
